\section{Evaluation Measures}
\label{sec:evaluation}

\subsection{Official metrics}

Let $\mathcal{L}$ contain the four task labels, $n_c$ be the number of gold instances with label $c$, $N$ the number of instances, and $F_c$ the label's F1 score. The official classification measure is \begin{equation} \operatorname{wF1}=\sum_{c\in\mathcal{L}}\frac{n_c}{N}F_c, \qquad F_c=\frac{2P_cR_c}{P_c+R_c}. \end{equation} Higher is better. Weighting lets common labels dominate the result. In our track-1 data, \NO{} is 16\% of the instances and receives the smallest weight, although it is the hardest label for the Assignment 1 models. We therefore report macro F1 beside the official score.

Let $\operatorname{Rev}$ leave \EQ{} unchanged and exchange
\FW{} with \BW{}. It is undefined for
\NO{}, so a \NO{} prediction on either side of a reversible pair counts as incompatible. For the set $\mathcal{P}$ of reversible source pairs
$(x,x^r)$, prediction function $f$, and gold function $y$, the consistency measures are
\begin{align}
  \operatorname{SoftCons}
  &=\frac{1}{|\mathcal{P}|}\sum_{(x,x^r)\in\mathcal{P}}
    \mathbb{I}\big[f(x)\!=\!\operatorname{Rev}(f(x^r))\big],\\
  \operatorname{HardCons}
  &=\frac{1}{|\mathcal{P}|}\sum_{(x,x^r)\in\mathcal{P}}
    \mathbb{I}\big[f(x)=y(x)\notag\\
  &\qquad\qquad\qquad\land\, f(x^r)=y(x^r)\big].
\end{align}
Both range from zero to one, and higher is better. \SoftCons{} rewards compatible predictions
whether or not they are correct. \HardCons{} rewards a pair only when both directions are
correct, so $\operatorname{HardCons}\leq\operatorname{SoftCons}$. Negative examples have no
reversed twin and affect weighted F1 only. All reported values are for English track 1, with
the consistency measures averaged over the 277 reversible dev pairs.

Symmetrized decoding makes the two predictions of a pair agree under the reversal rule with
\NO{} mapped to itself, so it raises \SoftCons{} by design. It does not fix it at 1.0. The
scorer counts a \NO{} answer on a reversible pair as incompatible, so the \SoftCons{} of
symmetrized predictions is one minus the share of reversible pairs answered \NO{}; the
preliminary run measured 0.906 and 0.986 for the two trained systems (\cref{sec:preliminary}).
Constant \EQ{} prediction does reach 1.000, while obtaining only 0.110 weighted F1 in
Assignment 1. \SoftCons{} therefore carries little discriminating information between systems
constrained this way; weighted F1 and \HardCons{} determine whether the constraint also
produces correct decisions.

\subsection{Secondary measures}

We report macro F1, accuracy, and precision, recall, and F1 for every label. A pair-outcome table separates pairs that are correct in both directions, consistent but wrong, inconsistent with one correct direction, and inconsistent with neither direction correct. We also count direction confusions, especially forward entailment predicted as backward entailment and the reverse. These diagnostics distinguish semantic errors from failures to respect direction. Finally, for every system we report the number of reversible pairs on which it predicts \NO{} in at least one order. The scorer counts each of them as incompatible whatever the other prediction is, so this count separates a \SoftCons{} shortfall caused by that label from one caused by contradictory directions.

\subsection{Subgroups}

The subgroup analysis uses three predeclared families. The length-cue slices mark whether
the longer phrase agrees with, misleads, or gives no information about the gold direction.
Phrase-overlap groups record whether both, one, or neither dev phrase occurs verbatim and
case-insensitively in train. Length groups use total words in the unordered pair: short has at
most five words, medium has six or seven, and long has more than seven. These cutoffs are the
nearest observed values to the 33rd and 67th percentiles of unique dev source pairs. We report
weighted F1, macro F1, and accuracy for every subgroup; \SoftCons{} and \HardCons{} are omitted
from sliced results because a filter can otherwise change or split their pair-level unit.

\begin{table}[t]
\centering
\small
\begin{tabular}{@{}lrrr@{}}
\toprule
Length group & $n$ & DeBERTa-13 & both+logreg \\
\midrule
Short ($\leq5$ words) & 306 & 0.843 & 0.647 \\
Medium (6--7 words) & 215 & 0.833 & 0.566 \\
Long ($>7$ words) & 139 & 0.765 & 0.540 \\
\bottomrule
\end{tabular}
\caption{Weighted F1 by length group. The aggregate hides weaker performance on long phrase
pairs. DeBERTa-13 is the encoder with seed 13. Counts and scores come from the fixed
Assignment 1 prediction files in the reliability run
(\texttt{experiments/2026-11-15-reliability-track1/}).}
\label{tab:length-subgroups}
\end{table}

The long-pair weakness occurs for all three DeBERTa seeds: seed 13 falls from 0.843 weighted
F1 on short pairs to 0.765 on long pairs. Phrase overlap is an even larger separator for seed
13, from 0.970 when both phrases occur in train to 0.757 when neither occurs. These are
descriptive slices, not additional hypothesis tests, and the ten-instance misleading-cue
group is too small for a stable effect estimate.
